% Part: incompleteness
% Chapter: arithmetization-syntax
% Section: introduction

\documentclass[../../../include/open-logic-section]{subfiles}

\begin{document}

\olfileid[fr]{inc}{art}{int}
\olsection{Introduction}

Pour relier la calculabilité et la logique, il nous faut pouvoir
parler des objets de la logique (symboles, termes, !!{formula}s,
!!{derivation}s), des opérations que l'on effectue sur eux, ainsi
que de leurs propriétés et de leurs relations, d'une manière qui
se prête au calcul. Nous pouvons le faire directement, en étudiant
les fonctions et les relations calculables portant sur les symboles,
les suites de symboles et les autres objets construits à partir
d'eux. Comme les objets de la syntaxe logique sont tous finis et
formés à partir d'ensembles !!{enumerable}s de symboles, certains
modèles de calcul permettent cette approche. Mais d'autres modèles,
comme celui des fonctions récursives, ne portent que sur les nombres
et sur les fonctions et les relations entre nombres. De plus, nous
voulons finalement pouvoir traiter de la syntaxe au sein de
certaines théories, en particulier des théories formulées dans le
langage de l'arithmétique. Dans ces cas, il faut
\emph{arithmétiser} la syntaxe : représenter respectivement les
objets syntaxiques, les opérations sur ces objets et leurs
relations par des nombres, des fonctions arithmétiques et des
relations arithmétiques. L'idée, qui remonte à Leibniz, est
d'attribuer des nombres aux objets syntaxiques.

Il est relativement simple d'attribuer aux symboles des nombres qui
leur servent de « codes ». Certains symboles posent néanmoins un
petit problème : il existe, par exemple, une infinité de
!!{variable}s, et même une infinité de !!{function}s pour chaque
arité~$n$. On peut bien sûr attribuer systématiquement des codes
distincts aux symboles, de sorte que, par exemple, $\Obj v_2$ et
$\Obj v_3$ reçoivent des codes différents. Les suites de symboles,
comme les termes et les !!{formula}s, présentent une difficulté
plus grande. Mais si nous savons traiter arithmétiquement les suites
de nombres, par exemple en les codant à l'aide de puissances de
nombres premiers, nous pouvons étendre le codage des symboles
individuels aux suites de symboles, puis aux suites ou autres
agencements de !!{formula}s, comme les !!{derivation}s. Ce codage
étendu porte le nom de « numérotation de G\"odel ». Chaque terme,
chaque !!{formula} et chaque !!{derivation} reçoit un numéro de
G\"odel.

En codant les suites de symboles par les suites de leurs codes et
en choisissant un codage des suites que des fonctions calculables
permettent de manipuler, nous pouvons alors aussi manipuler les
numéros de G\"odel au moyen de fonctions calculables. En pratique,
toutes les fonctions pertinentes seront récursives primitives. Par
exemple, calculer la longueur d'une suite et calculer son $i$-ième
élément à partir de son code sont deux opérations récursives
primitives. Si le nombre qui code la suite est, par exemple, le
numéro de G\"odel d'!!a{formula}~$!A$, nous voyons aussitôt que la
longueur d'!!a{formula} et le code de son $i$-ième symbole peuvent
également être calculés à partir du numéro de G\"odel de~$!A$. Il
est un peu plus difficile de prouver que la propriété d'être le
numéro de G\"odel d'un terme bien formé ou d'!!a{derivation}
correcte est récursive primitive. C'est néanmoins possible parce
que les suites considérées (termes, !!{formula}s, !!{derivation}s)
sont définies inductivement.

Prenons pour exemple l'opération de substitution. Si $!A$ est une
formule, $x$ une variable et $t$ un terme libre pour~$x$ dans~$!A$, alors $\Subst{!A}{t}{x}$
résulte du remplacement de chaque occurrence libre de~$x$ dans~$!A$
par~$t$. Supposons maintenant que nous ayons attribué des numéros de
G\"odel à $!A$, $x$ et $t$, disons respectivement $k$, $l$ et $m$.
Le même codage attribue un numéro de G\"odel à $\Subst{!A}{t}{x}$,
disons~$n$. L'application qui associe $n$ à $k$, $l$ et $m$ est
l'analogue arithmétique de l'opération de substitution. Tandis que
cette opération envoie $!A$, $x$ et $t$ sur $\Subst{!A}{t}{x}$,
la fonction de substitution arithmétisée envoie les numéros de
G\"odel $k$, $l$ et $m$ sur le numéro de G\"odel~$n$. Nous
verrons que cette fonction est récursive primitive.

L'arithmétisation de la syntaxe n'a pas qu'un intérêt abstrait. Ce
fut d'abord une idée féconde, et nullement évidente, que des
langages comme celui de l'arithmétique, qui ne disposent pas au
départ de moyens pour « parler des » langages, puissent malgré tout
formaliser des propriétés complexes des expressions. Il suffit
alors d'un pas de plus pour se demander ce qu'une théorie formulée
dans un tel langage, comme l'arithmétique de Peano, peut
\emph{démontrer} à propos de son propre langage, notamment sur la
démontrabilité ou la vérité d'!!{sentence}s. On arrive ainsi aux
célèbres théorèmes de limitation de G\"odel, sur la
non-démontrabilité, et de Tarski, sur l'indéfinissabilité de la
vérité. Mais l'arithmétisation de la syntaxe est aussi essentielle
à la preuve de résultats importants en théorie de la calculabilité,
par exemple sur la puissance de calcul des théories ou sur les
rapports entre différents modèles de calcul. Elle sert en outre de
modèle pour l'arithmétisation d'autres objets et propriétés. On
peut ainsi arithmétiser les configurations et les calculs, par
exemple ceux des machines de Turing. Cela permet de simuler les
calculs d'un modèle, tel celui des machines de Turing, dans un autre,
tel celui des fonctions récursives.

\end{document}
